\section{Persistent witnesses: B\"uchi strategies and their dual certificates}
\label{sec:games}

\status{A finite total-arena B\"uchi solver and both certificate checkers are
implemented in Python. Generalized acceptance and source-level strategy
extraction are specified extensions, not executed features.}

\subsection{A strategy is a witness with obligations over its future}
Let $G=(V,V_0,V_1,E,B)$ be a finite directed arena. The ownership sets partition
$V$, every vertex has at least one successor, and $B\subseteq V$ is the set of
accepting vertices. Player 0 is the protagonist; player 1 is the antagonist.
A play is an infinite sequence $v_0v_1\ldots$ following $E$. Each owner chooses
the successor when its vertex is visited.

The positive question at an initial vertex $s$ is
\begin{equation}
 \exists\sigma\;\forall\pi\text{ consistent with }\sigma\text{ from }s,
 \quad \forall N\in\N\;\exists n\geq N:\ v_n\in B.
 \label{eq:buchi}
\end{equation}
The expression on the right is the B\"uchi condition, denoted $\GF B$.
A strategy may in general inspect the play history. On finite B\"uchi games,
positional strategies, depending only on the current vertex, suffice for both
players \cite{chatterjee-henzinger}. The shipped certificate language therefore
contains a finite policy table, not a program with unbounded history.

A single path through $B$ proves neither recurrence nor resistance to an
opponent. A witness synthesized separately for each finite sample proves
neither that the same policy works for all samples nor that it works forever.
The policy itself must be the persistent witness, and its behavioral law must
quantify over \emph{all} permitted future choices.

The input validator rejects deadlocks. Whether a deadlock should lose, win, or
be completed by a self-loop depends on the source semantics. Completing it
without an explicit convention changes the meaning of \eqref{eq:buchi}.
Likewise, any fairness restriction on an opponent must be represented as part
of the objective or the model; no fairness is assumed here.

\subsection{Positive certificate: decrease until acceptance, then reset}
A positive certificate consists of a region $W\subseteq V$ containing $s$, a
legal positional choice $\sigma(v)$ at every $v\in W\cap V_0$, and a natural
rank $\rho:W\to\N$. Define the allowed edges under this policy by
\[
 E_\sigma(v)=
 \begin{cases}
  \{\sigma(v)\},&v\in V_0,\\
  E(v),&v\in V_1.
 \end{cases}
\]
The checker verifies
\begin{align}
 &E_\sigma(v)\subseteq W &&(v\in W),\label{eq:win-closed}\\
 &\rho(v)>\rho(w) &&(v\in W\setminus B,\ w\in E_\sigma(v)).
 \label{eq:win-rank}
\end{align}
There is no decrease requirement when the \emph{source} vertex belongs to $B$.
At such a vertex the rank may rise again. Closure still applies there.

\begin{theorem}[Positive rank soundness]
\label{thm:buchi-rank}
An accepted positive certificate proves \eqref{eq:buchi} with its supplied
policy, extended arbitrarily outside $W$.
\end{theorem}
\begin{proof}
Closure keeps every consistent play in $W$. Suppose some such play visits $B$
only finitely often. Beyond its last accepting visit, every edge has a source
outside $B$, so \eqref{eq:win-rank} gives an infinite strictly decreasing
sequence of natural numbers. This is impossible. Equivalently, from any
nonaccepting position, at most its current rank many transitions can occur
before an accepting vertex is reached. Apply this argument after every finite
prefix to obtain the quantifiers in \eqref{eq:buchi}.
\end{proof}

This is not ordinary termination. For the two-vertex cycle $a\to b\to a$ with
$B=\{a\}$, take $\rho(a)=0$ and $\rho(b)=1$. The edge $b\to a$ decreases;
$a\to b$ resets the rank. Every run is infinite and accepting. A globally
strictly decreasing ranking would reject this correct recurrent behavior.
This is the specific difference from Forge's existing compilation of cyclic
proof obligations to well-founded descent.

There is a useful bounded response property hidden in the same certificate.
Let $H=\max_{v\in W}\rho(v)$. Starting at a nonaccepting vertex, acceptance
occurs within $H$ transitions. Starting at an accepting vertex, a later
accepting vertex occurs within $H+1$ transitions. The certificate thus gives a
coarse recurrence-gap bound as well as qualitative infinite recurrence.
This bound concerns graph steps; a source adapter with variable-duration steps
needs an additional cost or timing argument.

\subsection{Negative certificate: bound all future accepting visits}
A negative result must do more than exhibit one bad play. It must show that
\emph{no} protagonist strategy satisfies the objective. The dual witness is an
antagonist strategy $\tau$ that defeats every protagonist continuation.

Its certificate consists of $L\subseteq V$ containing $s$, a legal choice
$\tau(v)$ at every $v\in L\cap V_1$, and $\kappa:L\to\N$. Now the allowed
edges are
\[
 E_\tau(v)=
 \begin{cases}
  E(v),&v\in V_0,\\
  \{\tau(v)\},&v\in V_1.
 \end{cases}
\]
The local checks are
\begin{align}
 &E_\tau(v)\subseteq L,\label{eq:lose-closed}\\
 &\kappa(v)\geq\kappa(w)+\one_B(v)
       \qquad(w\in E_\tau(v)).\label{eq:lose-rank}
\end{align}
Unlike the positive rank, this rank never increases. Each departure from an
accepting vertex consumes at least one unit. Nonaccepting cycles of constant
rank are allowed, and are usually the intended counterbehavior.

\begin{theorem}[Negative rank soundness]
\label{thm:cobuchi-rank}
An accepted negative certificate forces at most $\kappa(s)$ accepting visits
on every infinite play consistent with $\tau$ from $s$. In particular,
\eqref{eq:buchi} is false.
\end{theorem}
\begin{proof}
For every finite prefix of length $N$, summing \eqref{eq:lose-rank} gives
\[
 \sum_{i=0}^{N-1}\one_B(v_i)
 \leq \kappa(v_0)-\kappa(v_N)
 \leq \kappa(s).
\]
Closure makes every rank in this expression defined. The bound holds for every
$N$, so infinitely many accepting visits are impossible. The certificate
checks every protagonist edge, hence the conclusion does not depend on a
particular protagonist policy.
\end{proof}

The distinction between a counterstrategy and a counterexample lasso is
operationally important. A lasso can refute one fixed candidate policy. It
cannot refute the existence of a different winning policy. The negative
certificate leaves all player-0 choices available and checks them universally.

\subsection{Search by repeated attractors}
For a total subarena $U$, the player-$i$ attractor of $T\subseteq U$ is the
least set containing $T$ and closed under the following additions: a
player-$i$ vertex is added when it has a successor already in the set; a vertex
owned by the other player is added when all its successors in $U$ are already
in the set. A layered construction records the round of addition and a
rank-decreasing choice at the controlled vertices.

The solver uses the classical repeated-attractor scheme discussed in
\cite{chatterjee-henzinger}, not that paper's faster specialized algorithm.
\begin{lstlisting}[caption={The implemented Buechi search, in semantic pseudocode.}]
U := all vertices
while U is nonempty:
    A := Attr_0(B intersect U) inside U
    D := U minus A
    if D is empty:
        return U as the winning region
    Z := Attr_1(D) inside U
    record player-1 choices that reach D and then avoid B
    remove Z from U
return the empty winning region
\end{lstlisting}

The maintained subarena is total: every remaining player-0 vertex has some
remaining successor, and every successor of a remaining player-1 vertex
remains. Thus the attractor computation never gains a spurious victory from an
empty successor set. This invariant also explains why the final attractor
policy remains winning in the original arena rather than only in an induced
graph with some inconvenient opponent edges removed.

When $D$ is nonempty, no vertex of $D$ is accepting. At a player-0 vertex of
$D$, every successor that remains in $U$ stays in $D$; at a player-1 vertex
there is a successor in $D$. The antagonist can therefore reach $D$ from $Z$
and avoid $B$ thereafter, unless the protagonist exits to an earlier removed
region. That earlier region is already losing. Layer numbers can decrease only
finitely often, so the assembled antagonist policy is co-B\"uchi winning.
At least one vertex is removed in each unsuccessful outer iteration.

When the final $U$ equals its attractor, the attractor levels supply the
positive rank. At nonaccepting player-0 vertices the recorded move descends; at
nonaccepting player-1 vertices all moves descend. At accepting player-0
vertices, choose any successor in $U$. No descent is required there.

The negative certificate intentionally does not expose the search's removal
layers as its authoritative argument. After fixing the assembled antagonist
policy on the losing region, the implementation solves the local longest-weight
problem
\begin{equation}
 \kappa(v)=\max_{w\in E_\tau(v)}
                  \bigl(\one_B(v)+\kappa(w)\bigr)
 \label{eq:counter-longest}
\end{equation}
by monotone iteration from zero. An accepting vertex cannot occur on a cycle
of this fixed-policy graph, or the policy would not force co-B\"uchi. Cycles
therefore have weight zero. Removing repeated zero-weight cycles from a walk
shows that its total accepting-departure weight has a finite maximum, and the
iteration stabilizes. The replay checker verifies only
\eqref{eq:lose-closed}--\eqref{eq:lose-rank}, not this optimization or the
attractor computations that preceded it.

\subsection{Completeness and computational costs}
\begin{proposition}[Completeness of the finite rank languages]
Every winning initial vertex in a finite total B\"uchi game has a positive
certificate of the preceding form; every losing initial vertex has a negative
certificate.
\end{proposition}
\begin{proof}
Choose positional winning strategies, whose existence follows from finite-game
memoryless determinacy. Under a winning protagonist policy, a reachable cycle
entirely outside $B$ would be a losing play. The subgraph on nonaccepting
vertices is therefore acyclic. Give an accepting vertex rank zero and a
nonaccepting vertex the maximum number of transitions needed to reach the
first accepting vertex in this policy graph. These ranks satisfy
\eqref{eq:win-rank}.

Under a winning antagonist policy, a reachable cycle containing an accepting
vertex would allow infinitely many accepting visits. No such cycle exists.
Give each vertex the maximum accepting-departure count of a finite walk from
it. Nonaccepting cycles can be erased without reducing this count; the maximum
is finite and satisfies \eqref{eq:lose-rank}. Restrict either construction to
the corresponding strategy-closed region.
\end{proof}

The implemented attractor uses repeated scans for transparency. With $n$
vertices and $m$ edges, one attractor may require $O(nm)$ work and there may be
$O(n)$ outer iterations, giving a coarse $O(n^2m)$ bound for this Python
version. It would be incorrect to label the implementation $O(nm)$ merely
because that bound is standard for the classical algorithm with better data
structures.

The direct production improvement is a predecessor queue with a counter of
not-yet-attracted successors for each opponent-owned vertex. Each edge is then
processed once per attractor, yielding $O(nm)$ for the outer algorithm. The
negative longest-weight iteration also admits an SCC-condensation dynamic
program. Neither optimization is required for soundness, and neither is
claimed to have been measured here. The rank checker itself costs $O(n+m)$
local tests, apart from integer bit costs and input validation.

\subsection{Examples that separate the semantics}
In the arena with edges $s\to s$, $s\to b$, $b\to b$ and $B=\{b\}$,
player 0 wins from $s$ when it owns $s$: choose $b$. Change only the owner of
$s$ to player 1 and the antagonist can remain at $s$ forever. The edge graph
has not changed, but the correct verdict has. Tests check both outcomes and
also reject the old policy after the ownership mutation.

In the arena $a\to b$, $b\to b$, with $B=\{a\}$ and start $a$, the target is
visited immediately. The B\"uchi objective is still false: there is only one
accepting visit. This is a regression against substituting ordinary
reachability for a temporal specification.

A final trap is to omit an inconvenient opponent edge during replay. The
positive checker reads all successors of every player-1 vertex from the
original problem. A certificate cannot shrink this list. Similarly, the
negative checker cannot remove a player-0 choice. These are semantic checks,
not optional certificate-format hygiene.

\subsection{A concrete next extension: several recurring obligations}
For a conjunction $\bigwedge_{j<k}\GF B_j$, a finite monitor can remember the
next acceptance set that must be visited. This is a specified extension, not
part of the shipped solver API. Use product states $(v,j)$, with $j\in\{0,
\ldots,k-1\}$, preserve the owner of $v$, and update on each departure from $v$:
\[
 j'=\begin{cases}(j+1)\bmod k,&v\in B_j,\\j,&v\notin B_j.\end{cases}
\]
Let the product accepting set consist of $(v,k-1)$ with $v\in B_{k-1}$.
A visit to this set completes a full round. Infinitely many completed rounds
are equivalent to infinitely many visits to every $B_j$; both directions follow
by following the next-required index. A positional product strategy becomes a
source strategy with $k$ memory states.

Crucially, ``the monitor index is zero'' is \emph{not} an adequate accepting
condition: an execution may remain at index zero forever without satisfying
$B_0$. The acceptance event must mean completion of a round. A future
implementation should test a system stuck before its first required event and
a system whose successful policy must alternate between two obligations.
This extension remains finite and precise; it does not silently claim support
for arbitrary fairness assumptions, partial observation, or general temporal
logic.
